\input _template

\setlanguage{EN}

\Title{Fun Algebra}

\MathcompsLink{fun-algebra}

\Author{Patrik Bak}

\sec Introduction

In this problem collection you will find my favorite problems -- ones with short, interesting, often tricky solutions that use various algebraic tricks.

\sec Problems

Without further ado, let's get to the problems. The key is to find a clever manipulation. It is hard to give a general recipe -- after all, harder problems really are largely about creativity.

\EnvId{uD63SrC0DjlBiSlD8WN1i-nested-operation-42-to-0}
\Problem{0}{}{
    Let
    $$
    m \circ n = \frac{mn+4}{m+n}.
    $$
    Determine
    $$
    (\dots((42 \circ 41) \circ 40) \circ \dots \circ 1) \circ 0.
    $$
}{
    The problem is full of seemingly random constants; for instance,~$42$ looks like it could be set artificially. If we had to evaluate such large numbers by hand, we would go crazy. Try the problem with smaller constants.
}{
    We have
    \begitems
    \i $1 \circ 0 = 4$
    \i $(2 \circ 1) \circ 0 = 2$
    \i $((3 \circ 2) \circ 1) \circ 0 = 2$
    \i $(((4 \circ 3) \circ 2) \circ 1) \circ 0 = 2$
    \enditems
    Isn't it suspicious that after a while it always comes out to 2?
}{
    \Answer{$2$.}

    Notice a special property of the number $2$ with respect to the operation $\circ$. For any $m \ne -2$ we have
    $$
    m \circ 2 = \frac{2m+4}{m+2} = \frac{2(m+2)}{m+2} = 2.
    $$
    We evaluate the expression from the problem step by step from the left, with the numbers entering the operation in the order $42, 41, \dots, 2, 1, 0$. The moment we combine the running result with the number $2$, the value of the whole expression turns into $2$. Every subsequent operation has the form $2 \circ k$, which is again $2$. The answer is therefore $2$.
}

\EnvId{r_xTFMkLNh1OKZqL4f6KP-reciprocals-of-pairwise-sums-identity}
\Problem{0}{}{
    The numbers $a,b,c$ are nonzero real numbers with sum $0$. Prove that
    $$
    \frac{1}{a+b}+\frac{1}{b+c}+\frac{1}{c+a}=\frac{a^2+b^2+c^2}{2abc}.
    $$
}{
    We want to cleverly use the condition $a+b+c=0$ to get rid of the complicated fractions on the left.
}{
    The key is to realize that $a+b=-c$, so the first fraction is actually equal to $-\frac 1c$, and similarly for the rest.
}{
    After simplifying, it suffices to show
    $$
    -\left(\frac 1a + \frac 1b + \frac 1c\right) = \frac{a^2+b^2+c^2}{2abc},
    $$
    which after expanding becomes
    $$
    2ab+2bc+2ca=a^2+b^2+c^2.
    $$
    Now we need to use the condition $a+b+c=0$ again. What could we do with it to obtain the expressions we are supposed to prove something about?
}{
    From the condition $a+b+c=0$ it follows that $a+b=-c$, $b+c=-a$, and~$c+a=-b$. We therefore rewrite the left-hand side of the identity to be proved as
    $$
    -\frac{1}{c} - \frac{1}{a} - \frac{1}{b} = -\frac{ab+bc+ca}{abc}.
    $$
    At the same time, squaring the equation $a+b+c=0$ gives
    $$
    0 = (a+b+c)^2 = a^2+b^2+c^2 + 2(ab+bc+ca),
    $$
    from which we express $ab+bc+ca = -\frac{1}{2}(a^2+b^2+c^2)$. Substituting this expression into the rewritten left-hand side, we immediately obtain
    $$
    -\frac{-\frac{1}{2}(a^2+b^2+c^2)}{abc} = \frac{a^2+b^2+c^2}{2abc},
    $$
    which is what we wanted to prove.
}

\EnvId{T30uF67thgCpvHWXzfoqr-values-of-cyclic-fraction-sum}
\Problem{0}{}{
    The numbers $a,b,c$ are positive reals such that $ab+bc+ca=1$. Determine all possible values of the expression
    $$
    \frac{a(b^2+1)}{a+b}+
    \frac{b(c^2+1)}{b+c}+
    \frac{c(a^2+1)}{c+a}.
    $$
}{
    The trick is to simplify each fraction individually by using the condition $ab+bc+ca=1$.
}{
    The expression we are studying is full of letters, and there is one number in it, $1$. How about getting rid of it?
}{
    \Answer{$2$.}

    Let us rewrite the first fraction by substituting $1 = ab+bc+ca$ into the numerator:
    $$
    \gather
    \frac{a(b^2+ab+bc+ca)}{a+b} = \frac{a(b(a+b)+c(a+b))}{a+b} =\\= \frac{a(a+b)(b+c)}{a+b} = ab+ac.
    \endgather
    $$
    Cyclically, for the second and~third fractions we obtain $bc+ab$ and $ca+bc$. Adding all three rewritten expressions, we get
    $$
    (ab+ac) + (bc+ab) + (ca+bc) = 2(ab+bc+ca).
    $$
    Since $ab+bc+ca=1$, the value of the expression is $2$.

    We still need to show that this value is attainable -- but clearly there exist positive numbers $a,b,c$ for which $ab+bc+ca=1$, for example $a=b=c=\sqrt{3}/3$.
}

\EnvId{iD6Yxw2GEuWjliHpq83ij-fraction-sum-with-product-one}
\Problem{0}{}{
    The numbers $a,b,c$ are positive reals with product 1. Determine all possible values of the expression
    $$
    \frac{1}{1+a+ab}+
    \frac{1}{1+b+bc}+
    \frac{1}{1+c+ca}.
    $$
}{
    The key is to perform some clever manipulation of certain fractions that lets us use the condition $abc=1$. What is a common thing we know how to do with fractions?
}{
    A simple way to use $abc=1$ in, say, the first fraction is to multiply it by~$c/c$; then we have
    $$
    \frac{1}{c+ac+abc}=\frac{1}{c+ac+1}.
    $$
    This fraction happens, by pure chance, to have the same denominator as the last fraction.
}{
    \Answer{$1$.}

    We multiply the first fraction by $c/c$ and the second by $ac/ac$. Using $abc=1$, we thereby rewrite their denominators into a form identical with the third fraction:
    $$
    \align
    \frac{1}{1+a+ab} &= \frac{c}{c+ac+abc} = \frac{c}{c+ac+1}, \\
    \frac{1}{1+b+bc} &= \frac{ac}{ac+abc+ab c^2} = \frac{ac}{ac+1+c}.
    \endalign
    $$
    Adding all three fractions, we get
    $$
    \frac{c}{1+c+ac} + \frac{ac}{1+c+ac} + \frac{1}{1+c+ac} = \frac{c+ac+1}{1+c+ac} = 1.
    $$
    The only possible value of the expression is $1$. The example $a=b=c=1$ shows that this value is attainable.
}

\EnvId{ZaI7YVI7CM2Q9THxNgNIu-nonnegative-product-of-three-factors}
\Problem{0}{}{
    The real numbers $a, b, c$ have sum 1. Prove that the number
    $$
    (a + bc)(b + ca)(c + ab)
    $$
    is nonnegative.
}{
    The trick is to rewrite each parenthesis individually by using the condition $a+b+c=1$.
}{
    A good indicator of how to rewrite each parenthesis is the fact that in each one we have $a+bc$, where the first term has degree~1 and~the second degree~2 -- how about somehow making them both degree~2?
}{
    We use the condition $a+b+c=1$ to rewrite the expressions in the parentheses. For the first parenthesis we have:
    $$
    \gather
    a + bc = a(a+b+c) + bc = a^2+ab+ac+bc =\\= a(a+b)+c(a+b) = (a+c)(a+b).
    \endgather
    $$
    Analogously we derive $b+ca = (a+b)(b+c)$ and~$c+ab = (a+c)(b+c)$. Multiplying all three expressions we get
    $$
    (a+bc)(b+ca)(c+ab) = (a+b)^2(b+c)^2(c+a)^2.
    $$
    The result is a product of squares of real numbers, and~is therefore nonnegative.
}

\EnvId{2aaFNpMFewtsuCcfuo_fD-reciprocal-equation-for-odd-powers}
\Problem{0}{}{
    For nonzero real numbers $a,b,c$ we have
    $$
    \frac1a + \frac1b + \frac1c = \frac1{a+b+c}.
    $$
    Prove that for every odd positive integer $n$ we have
    $$
    \frac1{a^n} + \frac1{b^n} + \frac1{c^n} = \frac1{a^n+b^n+c^n}.
    $$
}{
    Forget what we're proving. When can the assumption even hold? Try to find all triples $a,b,c$ that satisfy this equality.
}{
    The key is to realize that after expanding the condition we get
    $$
    (a+b+c)(ab+bc+ca)-abc=0.
    $$
    This expression can surprisingly be factored into a product. From there it will be straightforward.
}{
    Let's rewrite the given condition by multiplying by the denominators:
    $$
    \gather
    \frac{ab+bc+ca}{abc} = \frac{1}{a+b+c},\\
    (a+b+c)(ab+bc+ca) - abc = 0.
    \endgather
    $$
    We rewrite the left-hand expression as a product by successively factoring out:
    $$
    \align
    (a+b+c)(ab+bc+ca) - abc &= \\ =
    (a+b+c)(ab + c(a+b)) - abc &= \\ =
    (a+b)ab + abc + (a+b)^2c + (a+b)c^2 - abc &= \\ =
    (a+b)(ab + ac + bc + c^2) &= \\ =
    (a+b)(a(b+c) + c(b+c))  &= \\ =
    (a+b)(b+c)(c+a).&
    \endalign
    $$
    So the equality holds if and only if at least two of the numbers are opposite (i.e. negatives of each other). Without loss of generality let $b=-a$.
    Then for every odd positive integer $n$ we have $b^n = (-a)^n = -a^n$.
    The left-hand side of the equality we're proving is
    $$
    \frac{1}{a^n} + \frac{1}{-a^n} + \frac{1}{c^n} = \frac{1}{c^n}.
    $$
    The right-hand side is
    $$
    \frac{1}{a^n - a^n + c^n} = \frac{1}{c^n}.
    $$
    The left-hand side equals the right-hand side, which completes the proof.
}

\EnvId{A_a1oWpvDYI-XyQHxUAwx-chain-of-reciprocal-equalities}
\Problem{0}{}{
    Pairwise distinct real numbers $a,b,c$ satisfy
    $$
    a + \frac1b = b + \frac1c = c + \frac1a.
    $$
    Prove that $|abc|=1$.
}{
    The problem tells us that
    $$
    \align
    a + \frac 1b &= b + \frac 1c, \\
    b + \frac 1c &= c + \frac 1a, \\
    c + \frac 1a &= a + \frac 1b. \\
    \endalign
    $$
    Realizing that this is a system of three equations is a good step toward making it natural to try combining the equations in various ways, as is usual when working with systems.
}{
    After a moment we easily see that adding and subtracting the equations leads nowhere. The key is to multiply these equations. To do that, though, we need to rearrange each equation a bit, so that something cancels when we multiply.
}{
    From the equality $a + \frac{1}{b} = b + \frac{1}{c}$ we express the difference $a-b$:
    $$
    a - b = \frac{1}{c} - \frac{1}{b} = \frac{b-c}{bc}.
    $$
    Cyclically, for the other pairs we get
    $$
    b - c = \frac{c-a}{ca} \quad\text{and}\quad c - a = \frac{a-b}{ab}.
    $$
    Since the numbers $a,b,c$ are pairwise distinct, the differences $a-b$, $b-c$ and~$c-a$ are nonzero. Multiplying these three equalities we get
    $$
    (a-b)(b-c)(c-a) = \frac{(b-c)(c-a)(a-b)}{(abc)^2}.
    $$
    Dividing by the nonzero expression $(a-b)(b-c)(c-a)$ we obtain $1 = \frac{1}{(abc)^2}$, hence $(abc)^2 = 1$, from which $|abc|=1$ follows.
}

\EnvId{ezUe8vAa8glYER7DWNTS6-system-of-square-differences}
\Problem{0}{}{
    For nonzero real numbers $a,b,c$ we have
    $$
    a^2-b^2 = bc \hbox{\quad and \quad} b^2-c^2 = ca.
    $$
    Prove that $a^2 - c^2 = ab$.
}{
    Adding the equations gives exactly the left-hand side of the equality we want to prove. So it suffices to prove something that looks seemingly simpler.
}{
    From the first hint we've already learned that adding the equations from the condition makes sense. That's not everything; we need something more. Besides adding the equations, there is for instance multiplying them. But first we often have to rearrange them -- ideally so that something cancels after multiplying.
}{
    Adding the given equations we get
    $$
    a^2-c^2 = bc+ca.
    $$
    So to prove the claim it suffices to show that $ab = bc+ca$.
    We rewrite the first given equation as $a^2 = b(b+c)$ and~the second as $(b-c)(b+c)=ca$. Multiplying these equations we have
    $$
    a^2(b-c)(b+c) = abc(b+c).
    $$
    If $b+c=0$, then from the first rewritten equation $a^2=0$, hence $a=0$, which contradicts the assumptions. So we can divide the equation by the nonzero expression $a(b+c)$, obtaining
    $$
    a(b-c) = bc,
    $$
    which is clearly equivalent to the relation $ab=bc+ca$ we set out to prove. This completes the proof.
}

\EnvId{OHiG7BBuj0eu_jJB1Xcmq-two-four-digit-factors}
\Problem{1}{}{
    Using only pen and paper (and your brain), find two four-digit numbers whose product equals $4^8 + 6^8 + 9^8$.
}{
    The problem hints that we should look for a factorization. Our expression is very unwieldy, though, so let us first introduce a suitable substitution of letters for the numbers to make it cleaner.
}{
    The forms $4=2^2$, $9=3^2$ and $6=2 \cdot 3$ suggest that the key substitution is $a=2$, $b=3$; the expression is then equal to
    $$
    a^{16}+(ab)^8+b^{16}.
    $$
    Can we factor this into a product?
}{
    The trick to factoring $a^{16}+(ab)^8+b^{16}$ is \textit{completing the square}; try completing $a^{16}+b^{16}$ to a square.
}{
    \Answer{The numbers $5521$ and~$8113$.}

    Let $a=2$, $b=3$. The expression from the statement is then equal to
    $$
    a^{16} + a^8 b^8 + b^{16}.
    $$
    We rewrite this expression by completing the square and then applying the formula for the difference of squares:
    $$
    \align
    a^{16}+a^8b^8+b^{16} &= (a^{16} + 2a^8b^8 + b^{16}) - a^8b^8 \\
    &= (a^8+b^8)^2 - (a^4b^4)^2 \\
    &= (a^8+b^8-a^4b^4)(a^8+b^8+a^4b^4).
    \endalign
    $$
    Substitute $a=2$ and~$b=3$. Then $a^4=16$, $b^4=81$, $a^4b^4 = 1296$.
    Further, $a^8=256$ and~$b^8=6561$, and $a^8+b^8 = 256+6561 = 6817$.
    The individual factors are therefore equal to:
    $$
    \align
    6817 - 1296 &= 5521, \\
    6817 + 1296 &= 8113.
    \endalign
    $$
    Both numbers we found are four-digit.
}


\EnvId{2ZxeP9uuZOf-7koaME3RI-quartic-equation-over-the-reals}
\Problem{1}{}{
    Over the reals, solve the equation
    $$
    x^4+4x^3+6x^2+4x=6.
    $$
}{
    The constants on the left-hand side are not random and should remind you of something.
}{
    That something the constants should remind you of is even taught in school and has a name.
}{
    \Answer{$x=-1\pm \root 4 \of 7$.}

    The left-hand side of the equation resembles the binomial theorem for 4~terms. Indeed,
    $$
    (x+1)^4 = x^4+4x^3+6x^2+4x+1.
    $$
    That is almost the left-hand side of the equation; we still need to add~1. The equation is therefore equivalent to the equation
    $$
    (x+1)^4 = 7.
    $$
    Over the reals the equation has two solutions $x=-1\pm \root 4 \of 7$.
}

\EnvId{Sno14U7w9Jv5l0joONjnr-product-of-three-quadratics-equation}
\Problem{2}{}{
    Over the reals, solve the equation
    $$
    (x^2 + 3x + 2)(x^2 - 2x - 1)(x^2 - 7x + 12) + 24 = 0.
    $$
}{
    The quadratic terms are not random. How about looking at each one separately -- can it be transformed?
}{
    The first step of the problem is to realize that two of our three quadratic terms can be factored into a product -- surprisingly it even comes out nicely, is that a coincidence?
}{
    What to do with four linear terms and~one quadratic? The trick is to build three quadratic terms out of them again.
}{
    Applying the advice from the previous hints, we should arrive at the equation
    $$
    (x^2 - 2x - 8)(x^2 - 2x - 3)(x^2 - 2x - 1) + 24=0.
    $$
    We can lower its degree with a suitable substitution. What should remain is an equation with integer roots that we can then solve by high-school methods.
}{
    \Answer{The solution set is $\{0, 2, 1 \pm \sqrt{6}, 1 \pm 2\sqrt{2}\}$.}

    We factor the first and~third parentheses into linear factors:
    $$
    (x+1)(x+2)(x^2-2x-1)(x-3)(x-4) + 24 = 0.
    $$
    By suitably regrouping the terms we get:
    $$
    \align
    (x+1)(x-3) &= x^2-2x-3 \\
    (x+2)(x-4) &= x^2-2x-8
    \endalign
    $$
    Substituting into the original equation we get
    $$
    (x^2-2x-3)(x^2-2x-8)(x^2-2x-1) + 24 = 0.
    $$
    We introduce the substitution $y = x^2-2x$. The equation turns into
    $$
    (y-3)(y-8)(y-1) + 24 = 0.
    $$
    After expanding and~simplifying we obtain
    $$
    \align
    y^3 - 12y^2 + 35y &= 0 \\
    y(y^2-12y+35) &= 0 \\
    y(y-5)(y-7) &= 0.
    \endalign
    $$
    For $y \in \{0, 5, 7\}$ we finish by solving the quadratic equations for $x$:
    \begitems
    \i $x^2-2x = 0$ has roots $0, 2$.
    \i $x^2-2x-5 = 0$ has roots $1 \pm \sqrt{6}$.
    \i $x^2-2x-7 = 0$ has roots $1 \pm 2\sqrt{2}$.
    \enditems
    The solution set is $\{0, 2, 1 \pm \sqrt{6}, 1 \pm 2\sqrt{2}\}$.
}


\EnvId{iwajcvpGTVmTHCDbGjT6a-squared-differences-and-cube-divisibility}
\Problem{2}{}{
    The integers $a,b,c$ satisfy
    $$
    (a-b)^2 + (b-c)^2 + (c-a)^2 = abc.
    $$
    Prove that $a+b+c+6 \mid a^3+b^3+c^3$.
}{
    The key is to recall an algebraic identity that somehow links some of the expressions this problem is about.
}{
    The key identity is the factorization of $a^3+b^3+c^3-3abc$ into a product. If you do not know it, try to derive this factorization. In the next hint we reveal possible tricks for how to find it.
}{
    To discover $a^3+b^3+c^3-3abc$ we can use these possible approaches:
    \begitems \style a
    \i Complete $a^3+b^3$ to a cube, that is, extract it from the expansion of $(a+b)^3$.
    \i Complete $a^3+b^3+c^3$ to a cube, that is, extract it from the expansion of $(a+b+c)^3=...$. From this we extract $a^3+b^3+c^3=...$. What remains are only the mixed terms such as $a^2b$, $ab^2$ and $abc$. These need to be suitably arranged so that we can factor something out.
    \enditems
}{
    We reveal the key identity
    $$
    a^3+b^3+c^3-3abc = (a+b+c)(a^2+b^2+c^2-ab-bc-ca).
    $$
    (The derivation hinted at in the previous hint can be found in the solution.)
}{
    How do we use the key identity? We need one more trick: realizing how the second factor of the factorization relates to the condition from the statement.
}{
    After expanding, the condition from the statement can be written as
    $$
    2(a^2+b^2+c^2-ab-bc-ca)=abc.
    $$
    As it happens, the second factor of our factorization has appeared here. Now it only remains to carefully put everything together.
}{
    First we derive the identity for the factorization of $a^3+b^3+c^3-3abc$. We complete the sum of the first two cubes to a cube of a sum using the relation
    $$
    a^3+b^3 = (a+b)^3 - 3ab(a+b).
    $$
    Then we apply the sum-of-cubes formula
    $$
    x^3+y^3=(x+y)(x^2-xy+y^2)
    $$
    to the terms $(a+b)^3$ and~$c^3$:
    $$
    \align
    a^3+b^3+c^3-3abc &=\\
    (a+b)^3 + c^3 - 3ab(a+b) - 3abc &=\\
    (a+b+c)((a+b)^2 - (a+b)c + c^2) - 3ab(a+b+c) &=\\
    (a+b+c)(a^2+2ab+b^2 - ac - bc + c^2 - 3ab) &= \\
    (a+b+c)(a^2+b^2+c^2-ab-bc-ca)&.
    \endalign
    $$
    Denote the second factor by $K = a^2+b^2+c^2-ab-bc-ca$. Since $a,b,c$ are integers, $K$ is also an integer.
    We expand the condition from the statement and rewrite it as
    $$
    (a^2-2ab+b^2) + (b^2-2bc+c^2) + (c^2-2ca+a^2) = abc,
    $$
    so that $2K=abc$. We substitute this relation into the derived identity:
    $$
    a^3+b^3+c^3 - 3(2K) = (a+b+c)K.
    $$
    Isolating the sum of the third powers we get
    $$
    a^3+b^3+c^3 = 6K + (a+b+c)K = K(a+b+c+6).
    $$
    The expression $a+b+c+6$ therefore divides $a^3+b^3+c^3$.

    \Remark We show one more alternative way of completing to a cube using the expansion of ${(a+b+c)^3}$. After expanding we have:
    $$
    \gather
    (a+b+c)^3 = a^3+b^3+c^3+ \\
    {}+ 3(a^2b+ab^2+b^2c+bc^2+c^2a+ca^2) + 6abc
    \endgather
    $$
    so that
    $$
    \gather
    a^3+b^3+c^3 - 3abc = (a+b+c)^3 \\
    {}- 3(a^2b+ab^2+b^2c+bc^2+c^2a+ca^2)- 9abc,
    \endgather
    $$
    From the last line we can factor out~$-3$; let us focus only on the inside of that parenthesis:
    $$
    \align
    (a^2b+ab^2+b^2c+bc^2+c^2a+ac^2) + 3abc &= \\
    a^2b+ab^2+b^2c+c^2b+c^2a+a^2c+3abc &= \\
    [a^2b+ab^2+abc]+[b^2c+c^2b+abc]+[c^2a+a^2c+abc] &=\\
    ab(a+b+c) + bc(b+c+a) + ca(c+a+b) &= \\
    (a+b+c)(ab+bc+ca).
    \endalign
    $$
    Now we just put it all together:
    $$
    \gather
    a^3+b^3+c^3 - 3abc = (a+b+c)^3 \\
    {}- 3(a^2b+ab^2+b^2c+bc^2+c^2a+ca^2)- 9abc = \\
    (a+b+c)^3 - 3(a+b+c)(ab+bc+ca) \\
    (a+b+c)((a+b+c)^2-3(ab+bc+ca)),
    \endgather
    $$
    which is an alternative way of writing the identity we are proving.
}

\DisplayHints

\DisplayProblemSolutions

\bye
